% ===== PRÉAMBULE CANONIQUE — à coller VERBATIM en tête de CHAQUE .tex =====
\documentclass[11pt,a4paper]{article}
\usepackage[utf8]{inputenc}\usepackage[T1]{fontenc}\usepackage{lmodern}
\usepackage[french]{babel}
\usepackage{amsmath,amssymb,mathtools}\usepackage{icomma}
\usepackage[version=4]{mhchem}          % \ce{...} : équations & formules
\usepackage{siunitx}\sisetup{locale=FR} % \qty{}{}, \si{}, \ang{}
\usepackage{chemfig}                    % \chemfig{...} : structures développées
\usepackage{xcolor}\usepackage{colortbl}
\usepackage{tikz}\usepackage{pgfplots}\pgfplotsset{compat=1.18}
\usepackage[most]{tcolorbox}\usepackage{enumitem}\usepackage{array,booktabs}
\usepackage[a4paper,margin=2.2cm]{geometry}
\usetikzlibrary{arrows.meta,calc,angles,quotes,positioning,patterns,backgrounds,shapes.geometric}
\definecolor{primary}{RGB}{222,49,99}\definecolor{primarydark}{RGB}{150,22,55}
\definecolor{defcol}{RGB}{0,114,178}\definecolor{methcol}{RGB}{52,92,148}
\definecolor{excol}{RGB}{0,135,90}\definecolor{attncol}{RGB}{200,40,55}
\tcbset{boxbase/.style={enhanced,breakable,boxrule=0.9pt,arc=2.5pt,
  left=3mm,right=3mm,top=2.3mm,bottom=2.3mm,fonttitle=\bfseries,coltitle=white}}
% --- boîtes TUTORIEL (noms EXACTS, ne pas renommer) ---
\newtcolorbox{cle}[1][]{boxbase,colback=defcol!6,colframe=defcol,
  title={\textsc{À retenir}\ifx\relax#1\relax\relax\else\textmd{ : \itshape #1}\fi}}
\newtcolorbox{methode}[1][]{boxbase,colback=methcol!7,colframe=methcol,sharp corners,
  title={\textsc{Méthode}\ifx\relax#1\relax\relax\else\textmd{ --- \itshape #1}\fi}}
\newtcolorbox{exemplebox}[1][]{boxbase,colback=excol!5,colframe=excol,
  title={\textsc{Exemple}\ifx\relax#1\relax\relax\else\textmd{ : \itshape #1}\fi}}
\newtcolorbox{piege}[1][]{boxbase,colback=attncol!6,colframe=attncol,
  title={\textsc{Piège}\ifx\relax#1\relax\relax\else\textmd{ : \itshape #1}\fi}}
\newtcolorbox{remarque}[1][]{boxbase,colback=black!4,colframe=black!45,
  title={\textsc{Remarque}\ifx\relax#1\relax\relax\else\textmd{ : \itshape #1}\fi}}
% --- boîtes EXERCICE / CORRIGÉ (noms EXACTS, auto-comptées) ---
\newtcolorbox[auto counter]{exo}[1][]{boxbase,colback=primary!4,colframe=primary,
  title={\textsc{Exercice}~\thetcbcounter\ifx\relax#1\relax\relax\else\textmd{ --- \itshape #1}\fi}}
\newtcolorbox[auto counter]{corrige}[1][]{boxbase,colback=excol!4,colframe=excol,
  title={\textsc{Corrigé}~\thetcbcounter\ifx\relax#1\relax\relax\else\textmd{ --- \itshape #1}\fi}}
\newcommand{\R}{\mathbb{R}}\newcommand{\N}{\mathbb{N}}\newcommand{\Z}{\mathbb{Z}}\newcommand{\ds}{\displaystyle}
% (un \usepackage{fancyhdr}+en-tête est possible ; il est ignoré côté web)
% =========================================================================
\begin{document}

\begin{exo}[coefficient contre indice]
On considère l'écriture « \ce{2 Na2CO3} ».
\begin{enumerate}[label=\alph*)]
  \item Combien de molécules (motifs) de carbonate de sodium représente-t-elle ?
  \item Combien d'atomes de sodium cela fait-il en tout ?
  \item Combien d'atomes au total ?
\end{enumerate}
\end{exo}

\begin{exo}[lire un hydrate]
On considère le composé \ce{CaCO3 . 2 H2O}.
\begin{enumerate}[label=\alph*)]
  \item Que signifie l'écriture « \ce{. 2 H2O} » ?
  \item Combien de molécules d'eau sont associées à chaque motif ?
  \item Quelle est la formule du composé anhydre correspondant ?
\end{enumerate}
\end{exo}

\begin{exo}[masse molaire d'un sel hydraté]
Le sulfate de cuivre(II) pentahydraté a pour formule \ce{CuSO4 . 5 H2O}. \textbf{Données :}
$M(\ce{Cu})=\qty{63,5}{\gram\per\mol}$, $M(\ce{S})=\qty{32}{\gram\per\mol}$,
$M(\ce{O})=\qty{16}{\gram\per\mol}$, $M(\ce{H})=\qty{1}{\gram\per\mol}$.
\begin{enumerate}[label=\alph*)]
  \item Retrouve la formule \ce{CuSO4} à partir des ions cuivre(II) et sulfate.
  \item Calcule la masse molaire du sulfate de cuivre \emph{anhydre} \ce{CuSO4}.
  \item Calcule la masse molaire de l'hydrate \ce{CuSO4 . 5 H2O}.
\end{enumerate}
\end{exo}

\begin{exo}[identifier un métal par un sulfate hydraté]
Un sulfate métallique heptahydraté a pour formule \ce{MSO4 . 7 H2O}. On en dissout \qty{0,20}{\mol}, soit
\qty{55,6}{\gram}. \textbf{Données :} $M(\ce{S})=\qty{32}{\gram\per\mol}$,
$M(\ce{O})=\qty{16}{\gram\per\mol}$, $M(\ce{H})=\qty{1}{\gram\per\mol}$.
\begin{enumerate}[label=\alph*)]
  \item Calcule la masse molaire du composé hydraté.
  \item Déduis-en la masse molaire de \ce{M}.
  \item Identifie le métal \ce{M}.
\end{enumerate}
\end{exo}

\begin{exo}[mise en situation : un soluté physiologique]
Pour préparer un soluté, on dissout des \emph{sels neutres} solides afin d'apporter les ions \ce{Na+},
\ce{K+}, \ce{Ca^2+} et \ce{Cl-}.
\begin{enumerate}[label=\alph*)]
  \item Écris le chlorure de sodium et le chlorure de potassium.
  \item Écris le chlorure de calcium, et indique combien d'ions \ce{Cl-} il libère par motif.
  \item Pourquoi ces sels sont-ils neutres, alors qu'ils sont formés d'ions chargés ?
\end{enumerate}
\end{exo}

\end{document}
